%FILL THESE IN
\def\mytitle{Coursework Report}
\def\mykeywords{Computer, Graphics, Napier, Pyramid, OpenGL, Edinburgh, Desert, Egypt}
\def\myauthor{Bradley Jones}
\def\contact{40167459@live.napier.ac.uk}
\def\mymodule{Computer Graphics (SET08116)}
%YOU DON'T NEED TO TOUCH ANYTHING BELOW
\documentclass[10pt, a4paper]{article}
\usepackage[a4paper,outer=1.5cm,inner=1.5cm,top=1.75cm,bottom=1.5cm]{geometry}
\twocolumn
\usepackage{graphicx}
\graphicspath{{./images/}}
%colour our links, remove weird boxes
\usepackage[colorlinks,linkcolor={black},citecolor={blue!80!black},urlcolor={blue!80!black}]{hyperref}
%Stop indentation on new paragraphs
\usepackage[parfill]{parskip}
%% all this is for Arial
\usepackage[english]{babel}
\usepackage[T1]{fontenc}
\usepackage{uarial}
\renewcommand{\familydefault}{\sfdefault}
%Napier logo top right
\usepackage{watermark}
%Lorem Ipusm dolor please don't leave any in you final repot ;)
\usepackage{lipsum}
\usepackage{xcolor}
\usepackage{listings}
%give us the Capital H that we all know and love
\usepackage{float}
%tone down the linespacing after section titles
\usepackage{titlesec}
%Cool maths printing
\usepackage{amsmath}
%PseudoCode
\usepackage{algorithm2e}

\titlespacing{\subsection}{0pt}{\parskip}{-3pt}
\titlespacing{\subsubsection}{0pt}{\parskip}{-\parskip}
\titlespacing{\paragraph}{0pt}{\parskip}{\parskip}
\newcommand{\figuremacro}[5]{
    \begin{figure}[#1]
        \centering
        \includegraphics[width=#5\columnwidth]{#2}
        \caption[#3]{\textbf{#3}#4}
        \label{fig:#2}
    \end{figure}
}

\lstset{
	escapeinside={/*@}{@*/}, language=C++,
	basicstyle=\fontsize{8.5}{12}\selectfont,
	numbers=left,numbersep=2pt,xleftmargin=2pt,frame=tb,
    columns=fullflexible,showstringspaces=false,tabsize=4,
    keepspaces=true,showtabs=false,showspaces=false,
    backgroundcolor=\color{white}, morekeywords={inline,public,
    class,private,protected,struct},captionpos=t,lineskip=-0.4em,
	aboveskip=10pt, extendedchars=true, breaklines=true,
	prebreak = \raisebox{0ex}[0ex][0ex]{\ensuremath{\hookleftarrow}},
	keywordstyle=\color[rgb]{0,0,1},
	commentstyle=\color[rgb]{0.133,0.545,0.133},
	stringstyle=\color[rgb]{0.627,0.126,0.941}
}

\thiswatermark{\centering \put(336.5,-38.0){\includegraphics[scale=0.8]{logo}} }
\title{\mytitle}
\author{\myauthor\hspace{1em}\\\contact\\Edinburgh Napier University\hspace{0.5em}-\hspace{0.5em}\mymodule}
\date{}
\hypersetup{pdfauthor=\myauthor,pdftitle=\mytitle,pdfkeywords=\mykeywords}
\sloppy
\begin{document}
	\maketitle
	\begin{abstract}
		The purpose of this project was to develop a 3D scene using OpenGL and C++ containing multiple pieces of three dimensional geometry, texturing effects, lighting effects and some form of post processing effects. Several different methods and techniques were used to achieve a satisfactory end result. This report will detail all the technical aspects of the project, such as the code that was used to achieve some of the effects in question.
		
		
	\end{abstract}
    
	\textbf{Keywords -- }{\mykeywords}
    %START FROM HERE
	\section{Introduction}
	My scene known as "Pyramids" was chosen as a coursework task because of the opportunities it allowed in creating a visually distinct scene, as well as personal interest in ancient Egypt. Currently, the coursework was focused on adding the appropriate objects and textures to create a scene reminiscent of a ruin near the Great Pyramids of Giza.
	
	
	 \figuremacro{h}{pyramids}{Pyramids of Giza}{ - Photograph of the real world pyramids, used as reference material.}{1.0}
	\{Wikipedia\}.
	
    \section{1 Overview}
    The project was created in Visual Studio 2015 and has been focused on the creation of objects, texturing them and the addition of light sources. 
    It is hoped that lights and normal mapping will be implemented in the near future.
    
    During the development of the current scene many other ideas were taken into consideration, for example a woodland scene, however in the end a desert scene inspired by the pyramids seemed like it would be the most unique setting for a OpenGL project.
    
    
    \section{Methods/Techniques}
    
    No external libraries out-side of those provided were used in the
    implementation of the application.
    
    The project uses a great variety of different effects in an effort to create a aesthetically pleasing scene: 
    \begin{itemize}
    	\item \dots{}
    	\begin{itemize}
    		\item Textures
    		\item Point Lights
    	\end{itemize}
    \end{itemize}
    
    
    One technique that was used to great effect was texturing. This is a critical technique as poor textures can severely hamper the realism of a 3D project, while good textures can greatly enhance a scene. An example of this is the use of a appropriate pyramid brick texture on the objects. As it makes it look like the bricks are getting smaller and smaller as they reach the apex of the structure, as well as giving off the impression of depth in the object without the use of normal mapping.
    
     \figuremacro{h}{wip}{Example of Texture Use}{ - A work in progress screen shot of the scene.}{1.0}
    
   Another technique that was used to great effect was the use of Lighting. Lighting is a critical part of the creation of a OpenGL scene, as it allows users to show of their objects and post processing techniques. Point light works in many regards to normal lighting. However, now our light source
   has a position and it has a area which it lights.
   
   
   \figuremacro{h}{point_light}{Example of Point Light Use}{ - Another work in progress screen shot of the scene, one that demonstrates the early implementation of a point light}{1.0}
   
   Illustrated below is a short section of code from the .cpp file, showing this
    
 \begin{lstlisting}[caption = The section of code used for creating the point light, taken from the main.cpp file] 

// Set lighting values, Position (-25, 10, -10)
light.set_position(vec3(-10.0f, 20.0f, -30.0f));
// Light colour white
light.set_light_colour(vec4(1.0f, 1.0f, 1.0f, 1.0f));
// Set range to 20
light.set_range(2000);
// Load in shaders
eff.add_shader("shaders/point.vert", GL_VERTEX_SHADER);
eff.add_shader("shaders/point.frag", GL_FRAGMENT_SHADER);

tex_eff.add_shader("27_Texturing_Shader/simple_texture.vert", GL_VERTEX_SHADER);
tex_eff.add_shader("71_Greyscale/greyscale.frag", GL_FRAGMENT_SHADER);

// Build effect
eff.build();
\end{lstlisting}

\section{Post Processing}

\figuremacro{h}{greyscale}{Example of the greyscale post processing effect}{ - A screen shot of greyscale being used.}{1.0}

\figuremacro{h}{masking}{Example of the masking post processing effect}{ - A screen shot of masking being used.}{1.0}

As shown in the screenshots, two post-process effects have been implemented. The
first one, greyscale, works by modifying the colour values
of the existing scene in a fragment shader. The second effect, masking is achieved by overplaying a checkered texture onto the scene. Both of these effects can be toggled on and off by using the 1/2 and 3/4 keys respectively.





	
\section{Conclusion}
The scene works mostly as expected. Several models have been loaded and textured. Point lighting and various  other shading effects have been implemented to create an interesting scene, however, some other effects have been removed from the scene as they were not completed or were clashing with existing aspects of the scene. Other effects like shadow mapping and skyboxes could have been implemented to improve the overall realism of the scene. 
\bibliographystyle{ieeetr}
\bibliography{references}		
\end{document}


